\subsection{Region Graph and B\'ezier Trajectory Model}
\label{sec:gcs-bezier-model}

The planner starts from the collision-free convex cover
$\mathcal{Q}=\{Q_i\}_{i=1}^{N_R}$ in
Eq.~\eqref{eq:convex-cover}, constructed by either ACD or VCC.  Each region
$Q_i$ is represented by a graph vertex $v_i$; we write $Q_{v_i}=Q_i$ when
indexing regions by vertices.  Thus $|V|=N_R$ and $|E|=N_E$.  The source
$\sigma$ and target $\tau$ are the designated vertices associated with the
regions containing $\mathbf{x}_{\mathrm{start}}$ and
$\mathbf{x}_{\mathrm{goal}}$, respectively.  A directed edge
$e=(u,v)$ represents a feasible transition from $Q_u$ to $Q_v$.  The discrete
problem selects an $\sigma$--$\tau$ path, while the continuous problem
determines a trajectory in its selected regions.

\begin{figure}[t]
  \centering
  \begin{tikzpicture}[scale=0.69, >=Latex, line cap=round, line join=round,
      every node/.style={font=\scriptsize}]
    \begin{scope}
      \draw[thick] (0,0) rectangle (6.1,3.6);
      \fill[black!40] (2.52,1.45) rectangle (3.20,3.10);
      \fill[black!40] (4.65,0.45) -- (5.65,0.45) -- (5.15,1.25) -- cycle;
      \filldraw[fill=blue!15, draw=blue!65!black, thick]
        (0.35,0.35) -- (2.15,0.35) -- (2.45,1.75) --
        (1.55,2.10) -- (0.45,1.45) -- cycle;
      \filldraw[fill=orange!17, draw=orange!75!black, thick]
        (1.55,0.75) -- (3.75,0.55) -- (4.45,1.80) --
        (3.85,2.75) -- (2.45,2.55) -- cycle;
      \filldraw[fill=green!17, draw=green!55!black, thick]
        (3.35,1.45) -- (5.90,1.55) -- (5.85,3.25) --
        (4.30,3.35) -- (3.35,2.60) -- cycle;
      \fill (0.90,0.78) circle (1.8pt);
      \fill (5.48,2.78) circle (1.8pt);
      \node[below left] at (0.90,0.78) {$\mathbf{x}_{\mathrm{start}}$};
      \node[above right] at (5.48,2.78) {$\mathbf{x}_{\mathrm{goal}}$};
      \node[blue!65!black] at (1.08,1.45) {$Q_\sigma$};
      \node[orange!75!black] at (3.17,1.30) {$Q_1$};
      \node[green!55!black] at (4.92,2.35) {$Q_\tau$};
      \node[below] at (3.05,0) {(a) convex cover};
    \end{scope}
    \draw[->, very thick] (6.62,1.80) -- (7.45,1.80);
    \begin{scope}[xshift=8.10cm]
      \node[circle, draw=blue!65!black, fill=blue!10, minimum size=6.4mm]
        (s) at (0.0,0.9) {$\sigma$};
      \node[circle, draw=orange!75!black, fill=orange!12, minimum size=6.4mm]
        (a) at (1.60,0.9) {$v_1$};
      \node[circle, draw=green!55!black, fill=green!10, minimum size=6.4mm]
        (t) at (3.50,0.9) {$\tau$};
      \draw[->] (s) -- (a);
      \draw[->] (a) -- (t);
      \node[above] at (1.75,2.10) {$G=(V,E)$};
      \node[below] at (1.75,0) {(b) induced region graph};
    \end{scope}
  \end{tikzpicture}
  \caption{From a convex free-space cover to a GCS.  Each cover region maps to
  one graph vertex, and each feasible common interface maps to a directed
  transition.  The start and goal lie in the designated source and target
  regions.}
  \label{fig:region-to-gcs}
\end{figure}

For a region $Q_v\subset\mathbb{R}^{n_c}$, the local trajectory is an
order-$d$ B\'ezier curve
\begin{equation}
  \mathbf{q}_v(s)
  =\sum_{k=0}^{d}B_{k,d}(s)\,\mathbf{r}_{v,k},
  \qquad s\in[0,1],
  \label{eq:framework-bezier-curve}
\end{equation}
where $\mathbf{r}_{v,k}\in\mathbb{R}^{n_c}$ are spatial control points and
$B_{k,d}$ are the Bernstein basis functions.  A second B\'ezier curve
$t_v(s)=\sum_{k=0}^{d}B_{k,d}(s)h_{v,k}$ parametrizes time.  The local decision
variable is consequently
\[
  x_v=(\mathbf{r}_{v,0},\ldots,\mathbf{r}_{v,d},
       h_{v,0},\ldots,h_{v,d}).
\]

The convex-hull property is the key reason for this representation: if every
control point belongs to the convex, collision-free region $Q_v$, the complete
curve belongs to $Q_v$.  In particular, when
$Q_v=\{\mathbf{q}\mid A_v\mathbf{q}\leq b_v\}$, continuous-time collision
avoidance is certified by the finite set of linear constraints
\begin{equation}
  A_v\mathbf{r}_{v,k}\leq b_v,\qquad k=0,\ldots,d.
  \label{eq:framework-containment}
\end{equation}
This converts continuous-time collision avoidance into finitely many convex
constraints~\cite{MarcucciEtAl2023MotionPlanning}.
The timing and velocity constraints are written in convex perspective form:
\begin{equation}
  \Delta h_{v,k}\geq\varepsilon,\qquad
  \Delta\mathbf{r}_{v,k}\in\Delta h_{v,k}D,\qquad k=0,\ldots,d-1,
  \label{eq:framework-kinematics}
\end{equation}
where $\Delta\mathbf{r}_{v,k}=\mathbf{r}_{v,k+1}-\mathbf{r}_{v,k}$,
$\Delta h_{v,k}=h_{v,k+1}-h_{v,k}$, and $\varepsilon>0$ prevents a degenerate
time scaling.  Here $D\subseteq\mathbb{R}^{n_c}$ is a prescribed convex set of
admissible configuration velocities, whose particular parametrization is set
by the planning model.  For each selected edge $(u,v)$,
position and time are joined through
\begin{equation}
  \mathbf{r}_{u,d}=\mathbf{r}_{v,0},
  \qquad h_{u,d}=h_{v,0}.
  \label{eq:framework-c0-continuity}
\end{equation}
Together with Eq.~\eqref{eq:framework-containment}, this places the shared
endpoint in $Q_u\cap Q_v$.  Higher-order continuity constraints can be included
in the edge model when a particular application requires them.

\subsection{Perspective GCS Formulation}
\label{sec:perspective-gcs}

Let $\mathcal{X}_v$ collect the local constraints in
Eqs.~\eqref{eq:framework-containment}--\eqref{eq:framework-kinematics}, and let
$\mathcal{X}_e$ collect the compatibility constraints for edge $e=(u,v)$.  For
the designated source and target regions, $\mathcal{X}_\sigma$ and
$\mathcal{X}_\tau$ additionally enforce
$\mathbf{r}_{\sigma,0}=\mathbf{x}_{\mathrm{start}}$ and
$\mathbf{r}_{\tau,d}=\mathbf{x}_{\mathrm{goal}}$, respectively.  GCS
uses a binary flow variable $y_e$ to indicate whether an edge is selected.
Rather than multiplying $y_e$ by the continuous variables directly, the
perspective formulation introduces edge-local copies $z_{e,u}$ and $z_{e,v}$.
For a convex edge cost $\ell_e$, its closed perspective is
\begin{equation}
  \widetilde{\ell}_e((z_{e,u},z_{e,v}),y_e)
  =y_e\,\ell_e\!\left(\frac{z_{e,u}}{y_e},
                       \frac{z_{e,v}}{y_e}\right)
  \quad \text{for } y_e>0,
  \label{eq:framework-perspective-cost}
\end{equation}
with the usual closed extension at $y_e=0$.  This construction preserves
convexity and switches the continuous variables and cost off when an edge is
inactive~\cite{MarcucciEtAl2023GCS}.

For a vertex $v$, let
$\mathcal{E}^{+}(v)=\{e=(v,w)\in E\}$ and
$\mathcal{E}^{-}(v)=\{e=(w,v)\in E\}$ denote its outgoing and incoming edge
sets, respectively.  For an edge $e=(u,v)$, $z_{e,u}$ is the copy of the local
trajectory variable at its tail $u$, and $z_{e,v}$ is the copy at its head $v$.
The resulting GCS--B\'ezier problem can be expressed compactly as
\begin{equation}
\begin{aligned}
  \min_{y,z}\quad &
    \sum_{e=(u,v)\in E}
    \widetilde{\ell}_e((z_{e,u},z_{e,v}),y_e) \\
  \text{s.t.}\quad &
    \sum_{e\in\mathcal{E}^{+}(v)}y_e-
    \sum_{e\in\mathcal{E}^{-}(v)}y_e=\delta_v,
    && \forall v\in V,\\
  &
    \sum_{e\in\mathcal{E}^{-}(v)}z_{e,v}=
    \sum_{e\in\mathcal{E}^{+}(v)}z_{e,u},
    && \forall v\in V\setminus\{\sigma,\tau\},\\
  &
    z_{e,u}\in y_e\mathcal{X}_u,\quad
    z_{e,v}\in y_e\mathcal{X}_v,\quad
    (z_{e,u},z_{e,v})\in y_e\mathcal{X}_e,
    && \forall e=(u,v)\in E,\\
  & y_e\in\{0,1\},
    && \forall e\in E .
\end{aligned}
\label{eq:framework-gcs-micp}
\end{equation}
Here, $\delta_\sigma=1$, $\delta_\tau=-1$, and $\delta_v=0$ otherwise.
The first constraint routes one unit of flow from $\sigma$ to $\tau$.  The second equates the
head copy of each incoming edge with the tail copy of each outgoing edge at the
same region, thereby enforcing state-flow consistency across local
trajectories.  Depending on the chosen convex costs and
sets, this model is a mixed-integer convex program or a mixed-integer
second-order cone program~\cite{MarcucciEtAl2023GCS,MarcucciEtAl2023MotionPlanning}.
It can combine convex penalties on path length, traversal time, and trajectory
smoothness.

\subsection{Convex Relaxation and Path Rounding}
\label{sec:gcs-relaxation-rounding}

Solving Eq.~\eqref{eq:framework-gcs-micp} to global optimality with
branch-and-bound can be expensive when the region graph is large.  The scalable
alternative used for the benchmark first relaxes the binary constraints to
$0\leq y_e\leq1$ and solves the resulting convex problem.  Its optimum
$C_{\mathrm{relax}}$ is a lower bound on the optimal mixed-integer cost
$C_{\mathrm{opt}}$.

The fractional flow is then converted to a valid discrete path by randomized
rounding.  At a current vertex $u$, an admissible outgoing edge is sampled
according to
\begin{equation}
  \Pr(e\mid u)
  =\frac{y_e^\star}
         {\sum_{e'\in\mathcal{E}^{+}(u)}y_{e'}^\star},
  \qquad e\in\mathcal{E}^{+}(u),
  \label{eq:framework-rounding}
\end{equation}
while excluding edges that would create a cycle.  Once a path reaches
$\tau$, its edge variables are fixed and the continuous trajectory variables
are re-optimized over that convex restriction.  The resulting feasible cost
$C_{\mathrm{round}}$ satisfies
\begin{equation}
  C_{\mathrm{relax}}\leq C_{\mathrm{opt}}\leq C_{\mathrm{round}}.
  \label{eq:framework-bounds}
\end{equation}
Accordingly, the reported relaxation--rounding gap
\begin{equation}
  g_{\mathrm{relax}}
  =\frac{C_{\mathrm{round}}-C_{\mathrm{relax}}}
         {C_{\mathrm{relax}}}
  \label{eq:framework-relaxation-gap}
\end{equation}
is a computable upper bound on the unknown relative suboptimality whenever
$C_{\mathrm{relax}}>0$.  A zero gap certifies optimality of the rounded
solution; otherwise, branch-and-bound is required for an exact global
certificate.  This distinction motivates reporting both graph size and
relaxation--rounding behavior in the experimental analysis.
